\documentclass[11pt]{article}
\usepackage[margin=1in]{geometry}
\usepackage{booktabs}
\usepackage{longtable}
\usepackage{array}
\usepackage{amsmath,amssymb}
\usepackage[hidelinks]{hyperref}
\usepackage{microtype}
\setlength{\parskip}{0.6em}
\setlength{\parindent}{0pt}

\begin{document}
\title{LSC 6.2.1 Codex Analysis}
\author{Open working package}
\date{\today}
\maketitle

\section*{Scope}
This report is an exploratory fit to the BEST, GALLEX and SAGE gallium source data.
It is not a claim of confirmed new physics. The aim is to test whether a compact
phenomenological ansatz can reproduce the main source-calibration ratios while
keeping the calibration trail readable.

\section*{Data Anchors}
The anchor set uses five independent constraints:
BEST inner, BEST outer, GALLEX combined reanalysis, SAGE $^51$Cr and SAGE $^37$Ar.
BEST ratio and the GALLEX Cr1/Cr2 split are kept as diagnostics.

\section*{Model}
The fit is performed in linearized form:
\[
-\ln R = \beta_L L^-1_{\mathrm eff} + \beta_S A_{\mathrm shape} +
\beta_\sigma \, \sigma_{\mathrm norm} + \beta_Z Z_{\mathrm BEST} + \beta_0.
\]
Here $L_{\mathrm eff}$ is a geometry proxy, $A_{\mathrm shape}$ is a second-moment
anisotropy proxy, $\sigma_{\mathrm norm}$ is the mean source-line cross section in
units of $10^{-45}\mathrm{cm}^2$, and $Z_{\mathrm BEST}$ splits the BEST inner and
outer zones.

\subsection*{Geometry Proxy}
For this first-pass package:
\[
L_{\mathrm eff} =
\begin{cases}
4R/3 & \text{sphere} \\
(R + H/2)/2 & \text{cylindrical or tank geometry}
\end{cases}
\]
\[
A_{\mathrm shape} = \frac{\left| \langle z^2 \rangle - \langle x^2 \rangle \right|}
{\langle z^2 \rangle + 2 \langle x^2 \rangle},
\quad
\langle x^2 \rangle = R^2/4,\quad \langle z^2 \rangle = H^2/12.
\]

\section*{Fitted Coefficients}
\begin{tabular}{lr}
\toprule
coefficient & value \\
\midrule
beta\_length & -0.39830548 \\beta\_shape & -0.71268732 \\beta\_source & 0.14607627 \\beta\_zone & 0.01969744 \\beta\_intercept & -0.13689794 \\
\bottomrule
\end{tabular}

\section*{Anchor Predictions}
\begin{longtable}{lrrr}
\toprule
case & observed & predicted & residual \\
\midrule
\endfirsthead
\toprule
case & observed & predicted & residual \\
\midrule
\endhead
BEST\_inner & 0.790000 & 0.790000 & +0.000000 \\BEST\_outer & 0.770000 & 0.770000 & -0.000000 \\GALLEX\_combined & 0.882000 & 0.882000 & +0.000000 \\SAGE\_51Cr & 0.950000 & 0.950000 & +0.000000 \\SAGE\_37Ar & 0.790000 & 0.790000 & +0.000000 \\
\bottomrule
\end{longtable}

\section*{Validation Predictions}
\begin{longtable}{lrrr}
\toprule
case & observed & predicted & residual \\
\midrule
\endfirsthead
\toprule
case & observed & predicted & residual \\
\midrule
\endhead
BEST\_ratio & 0.970000 & 0.974684 & +0.004684 \\GALLEX\_Cr1 & 0.953000 & 0.882000 & -0.071000 \\GALLEX\_Cr2 & 0.812000 & 0.882000 & +0.070000 \\
\bottomrule
\end{longtable}

\section*{Leave-One-Out Stability}
\begin{longtable}{lrrr}
\toprule
omitted case & observed & predicted & residual \\
\midrule
\endfirsthead
\toprule
omitted case & observed & predicted & residual \\
\midrule
\endhead
BEST\_inner & 0.790000 & 0.965858 & +0.175858 \\BEST\_outer & 0.770000 & 0.941584 & +0.171584 \\GALLEX\_combined & 0.882000 & 0.583630 & -0.298370 \\SAGE\_51Cr & 0.950000 & 0.789343 & -0.160657 \\SAGE\_37Ar & 0.790000 & 0.817808 & +0.027808 \\
\bottomrule
\end{longtable}

\section*{Interpretation}
The exact anchor fit means the main question is not whether the algebra can absorb the
anchors. It is whether the same coefficients survive perturbation. The leave-one-out run
shows that the package is still sensitive to the selected anchor set, especially around
GALLEX. That is useful: it tells us where the next 6.2.1 iteration needs extra nuisance
structure or additional event-level data.

\section*{Source Notes}
\begin{itemize}
\item BEST: \href{https://doi.org/10.1103/PhysRevLett.128.232501}{Phys. Rev. Lett. 128, 232501 (2022)}
\item GALLEX final and reanalysis: \href{https://doi.org/10.1016/S0370-2693(97)01562-1}{Phys. Lett. B 420, 114-126 (1998)} and \href{https://doi.org/10.1016/j.physletb.2010.01.030}{Phys. Lett. B 685, 47-54 (2010)}
\item SAGE Cr-51: \href{https://doi.org/10.1103/PhysRevC.59.2246}{Phys. Rev. C 59, 2246 (1999)}
\item SAGE Ar-37: \href{https://doi.org/10.1103/PhysRevC.73.045805}{Phys. Rev. C 73, 045805 (2006)}
\item Gallium source cross sections: \href{https://doi.org/10.1103/PhysRevC.56.3391}{Phys. Rev. C 56, 3391 (1997)}
\end{itemize}

\end{document}
